\chapter*{译者说明}
\addcontentsline{toc}{chapter}{译者说明}
\markboth{译者说明}{译者说明}

\textbf{底本。}本书译自 J. Bernstein 与 S. Gelbart 编《An Introduction to the Langlands Program》（Birkhäuser, 2004）。原书 12 章来自 2001 年 3 月希伯来大学高等研究院讲习班，六位作者分写：Kowalski（1–3 章）、de Shalit（4–5 章）、Kudla（6–7 章）、Bump（8 章）、Cogdell（9–11 章）、Gaitsgory（12 章）。

\textbf{页次。}所用底本 PDF 与印刷本相比缺 6 页，故 PDF 页与印刷页之差从 8 逐段递减到 2（第 12 章为 2）。译文引用页码一律用印刷页码；各章 PDF 页与印刷页的对应见工程文件 \texttt{work/PAGE\_MAP.md}。

\textbf{编排。}译文保留原书的章、节、编号与定理体例，不合并、不删节；原书的脚注照译，各章末参考文献照录。唯第 9--11 章书目条目繁多（分别 49、60、36 条），按体例以「见原书页码」指代，未逐条迻录。第 12 章（几何朗兰兹）此前已单独成册，本译本将其并入全书，并统一了术语与记号。

\textbf{术语与译注。}全书术语先冻结后动笔，术语表见书末。译注分三类：易错预警、含糊点拨、背景直觉；每条译注末句给出依据，A 为原书或作者论文、经典教材，B 为 nLab、Stacks Project 等白名单工具书（均与原书交叉核对），I 为本译丛既有译法。

\textbf{公式。}一切公式逐符号对照原书页面图像核对，不以文字提取件为准。若原书有排印笔误而影响理解，保留原样并加译注说明。

\textbf{分工。}全书各章由数译逐章译校，体例统一，文责自负。\hfill 数译，2026 年
